\documentclass[a4paper,10pt]{article}
\usepackage{eumat}

\begin{document}
\begin{eulernotebook}
\eulerheading{Statistika dengan Euler Math Toolbox (EMT)}
\begin{eulercomment}
Catatan ini membahas penggunaan Euler Math Toolbox (EMT) untuk
mengolah dan menganalisis data statistika. Materi yang dibahas
meliputi penyimpanan data, simulasi, pembacaan file, statistika
deskriptif, penyajian grafik, dan pengujian hipotesis. Setiap bagian
dilengkapi dengan penjelasan konsep serta contoh penggunaan perintah
EMT. Penyusunan materi mengacu pada notebook EMT4Statistika dari dosen
dengan tetap mempertahankan urutan pembahasannya. Pada bagian akhir,
terdapat analisis data dan lima soal latihan beserta penyelesaiannya.

\begin{eulercomment}
\eulerheading{Cara Menggunakan Notebook}
\begin{eulercomment}
Jalankan perintah EMT yang diawali tanda \textgreater{} secara berurutan, mulai
dari bagian awal. Beberapa perintah membutuhkan file data seperti
table.dat, table1.dat, excel-data.csv, dan sample.csv, sehingga file
tersebut perlu tersedia di folder kerja EMT. Pada bagian simulasi,
hasil yang muncul dapat berbeda setiap kali perintah dijalankan.
Ulangi perintah tersebut untuk mengamati perubahan hasilnya.

\begin{eulercomment}
\eulerheading{Pengantar Statistika}
\begin{eulercomment}
Statistika digunakan untuk mengolah data agar informasi di dalamnya
lebih mudah dipahami dan dapat digunakan sebagai dasar pengambilan
kesimpulan. Secara umum, pengolahan data dapat dilakukan dengan:\\
- Mengumpulkan data melalui survei, pengukuran, atau simulasi.\\
- Merangkum data menggunakan rata-rata, median, simpangan baku, dan
grafik.\\
- Mengamati pola atau hubungan yang terdapat dalam data.\\
- Melakukan pengujian statistika untuk menilai suatu dugaan
berdasarkan data yang diperoleh.

EMT membantu proses tersebut melalui berbagai perintah pengolahan dan
analisis data.

\end{eulercomment}
\eulersubheading{Beberapa istilah dasar yang digunakan dalam pembahasan:}
\begin{eulercomment}
- Rata-rata: jumlah seluruh nilai dibagi banyaknya data.\\
- Median: nilai tengah setelah data diurutkan.\\
- Simpangan baku: ukuran yang menunjukkan tingkat penyebaran data
terhadap rata-ratanya.\\
- Nilai-p: peluang memperoleh hasil yang sama ekstrem atau lebih
ekstrem daripada hasil pengamatan, jika hipotesis nol benar.

\begin{eulercomment}
\eulerheading{1. Data dan Statistika Deskriptif}
\begin{eulercomment}
\end{eulercomment}
\eulersubheading{Arti konsep}
\begin{eulercomment}
Data disimpan sebagai vektor (satu baris angka) atau matriks (baris =
pengamatan, kolom = variabel). Ukuran pemusatan menjawab "angka yang
mewakili data itu berapa", yaitu rata-rata dan median. Ukuran sebaran
menjawab "seberapa beragam datanya", yaitu simpangan baku (pembagi
n-1) dan kuartil.

\end{eulercomment}
\eulersubheading{Contoh dunia nyata}
\begin{eulercomment}
Gaji lima karyawan: 4, 4, 5, 5, dan 100 juta rupiah. Rata-ratanya 23,6
juta, padahal hampir semua orang bergaji 4 sampai 5 juta. Median 5
juta lebih jujur menggambarkan gaji yang umum. Jadi rata-rata mudah
terseret nilai ekstrem, sedangkan median tidak. Dalam praktik,
keduanya dilihat bersama.

\end{eulercomment}
\eulersubheading{Praktik di EMT}
\begin{eulercomment}
Tujuan: meringkas sepuluh hasil penimbangan bungkus gula yang
seharusnya 1000 gram.

Langkah 1. Simpan data sebagai vektor, lalu hitung median, rata-rata,
dan simpangan baku.
\end{eulercomment}
\begin{eulerprompt}
>M=[1000,1004,998,997,1002,1001,998,1004,998,997]; ...
>median(M), mean(M), dev(M),
\end{eulerprompt}
\begin{euleroutput}
  999
  999.9
  2.7264
\end{euleroutput}
\begin{eulercomment}
Hasil: median 999, rata-rata 999,9, simpangan baku 2,73. Data berkumpul di sekitar 1000
dan menyimpang rata-rata sekitar 2,73 gram. Pabrik memakai angka seperti ini untuk
menilai apakah mesin pengemas masih akurat.

Langkah 2. Kuartil Q1, Q2 (median), dan Q3 diambil dengan parameter p. Seperempat data
berada di bawah Q1 dan tiga perempat di bawah Q3.
\end{eulercomment}
\begin{eulerprompt}
>median(M,p=[0.25,0.5,0.75])
\end{eulerprompt}
\begin{euleroutput}
  [998,  999,  1001.8]
\end{euleroutput}
\begin{eulercomment}
Langkah 3. Diagram kotak (box-and-whiskers) merangkum nilai minimum, Q1, median, Q3, dan
maksimum. Pada data ini tidak ada pencilan (outlier).
\end{eulercomment}
\begin{eulerprompt}
>aspect(1.75); boxplot(M):
\end{eulerprompt}
\begin{eulercomment}
Langkah 4. Hitung peluang nilai lebih dari 1005, dengan anggapan data berasal dari
distribusi normal. Semua fungsi distribusi di Euler berakhiran "dis" dan menghitung
fungsi distribusi kumulatif (peluang X \textless{}= x):
\end{eulercomment}
\begin{eulerformula}
\[
\text{normaldis}(x,m,d)=\int_{-\infty}^x \frac{1}{d\sqrt{2\pi}}e^{-\frac{1}{2}\left(\frac{t-m}{d}\right)^2}\,dt
\]
\end{eulerformula}
\begin{eulercomment}
Hasilnya dicetak dalam persen dengan dua angka desimal memakai fungsi print.
\end{eulercomment}
\begin{eulerprompt}
>print((1-normaldis(1005,mean(M),dev(M)))*100,2,unit=" %")
\end{eulerprompt}
\begin{euleroutput}
        3.07 %
\end{euleroutput}
\begin{eulercomment}
Hasilnya 3,07 \%. Dari 100 bungkus, sekitar 3 bungkus diperkirakan lebih dari 1005 gram.
Catatan kritis: PDF dosen memakai gprint dan menampilkan galat "Function gprint not
found". Pada notebook (.en) perintahnya print, dan ini yang berjalan. Hasil ini hanya
sahih jika data memang mendekati normal, padahal sampelnya hanya sepuluh.

\end{eulercomment}
\eulersubheading{Distribusi frekuensi}
\begin{eulercomment}
Jika datanya ratusan, angka mentah sulit dibaca, jadi data dikelompokkan ke dalam
rentang (kelas). Misalkan banyak laki-laki pada tiap rentang tinggi badan adalah sebagai
berikut (r = batas kelas, v = frekuensi).
\end{eulercomment}
\begin{eulerprompt}
>r=155.5:4:187.5; v=[22,71,136,169,139,71,32,8];
\end{eulerprompt}
\begin{eulercomment}
Histogram dari data tersebut:
\end{eulercomment}
\begin{eulerprompt}
>plot2d(r,v,a=150,b=200,c=0,d=190,bar=1,style="\(\backslash\)/"):
\end{eulerprompt}
\begin{eulercomment}
Data itu dimasukkan ke dalam tabel. Tabel adalah cara menyimpan data statistika. Tabel ini
punya tiga kolom: batas bawah (BB), batas atas (BA), dan frekuensi (Frek). Tanda "\textbar{}"
memisahkan kolom, dan opsi labc pada writetable menentukan judul kolom.
\end{eulercomment}
\begin{eulerprompt}
>T:=r[1:8]' | r[2:9]' | v'; writetable(T,labc=["BB","BA","Frek"])
\end{eulerprompt}
\begin{euleroutput}
          BB        BA      Frek
       155.5     159.5        22
       159.5     163.5        71
       163.5     167.5       136
       167.5     171.5       169
       171.5     175.5       139
       175.5     179.5        71
       179.5     183.5        32
       183.5     187.5         8
\end{euleroutput}
\begin{eulercomment}
Untuk menghitung rata-rata dan sebaran, dibutuhkan titik tengah tiap kelas. Dengan dua
kolom pertama tabel:
\end{eulercomment}
\begin{eulerprompt}
>(T[,1]+T[,2])/2 // titik tengah tiap kelas
\end{eulerprompt}
\begin{euleroutput}
      157.5 
      161.5 
      165.5 
      169.5 
      173.5 
      177.5 
      181.5 
      185.5 
\end{euleroutput}
\begin{eulercomment}
Cara yang lebih singkat adalah melipat vektor r dengan [1/2,1/2]:
\end{eulercomment}
\begin{eulerprompt}
>M=fold(r,[0.5,0.5])
\end{eulerprompt}
\begin{euleroutput}
  [157.5,  161.5,  165.5,  169.5,  173.5,  177.5,  181.5,  185.5]
\end{euleroutput}
\begin{eulercomment}
Rata-rata dan simpangan baku sampel dengan frekuensi v:
\end{eulercomment}
\begin{eulerprompt}
>\{m,d\}=meandev(M,v); m, d,
\end{eulerprompt}
\begin{euleroutput}
  169.9
  5.9891
\end{euleroutput}
\begin{eulercomment}
Hasilnya rata-rata 169,90 dan simpangan baku 5,99 dari 648 orang.

Kurva kerapatan normal dengan rata-rata m dan simpangan baku d memakai rumus
\end{eulercomment}
\begin{eulerformula}
\[
y=\frac{1}{d\sqrt{2\pi}}e^{-\frac{(x-m)^2}{2d^2}}
\]
\end{eulerformula}
\begin{eulercomment}
Nilainya antara 0 dan 1, sehingga untuk ditumpuk pada histogram perlu dikali jumlah data
dan lebar kelas (4).
\end{eulercomment}
\begin{eulerprompt}
>plot2d("qnormal(x,m,d)*sum(v)*4", ...
>  xmin=min(r),xmax=max(r),thickness=3,add=1):
\end{eulerprompt}
\begin{eulercomment}
Kurva mengikuti bentuk histogram, jadi tinggi badan kira-kira berdistribusi normal.
Contoh nyata: toko pakaian memakai pola ini untuk menentukan berapa banyak ukuran S, M,
dan L yang perlu diproduksi.

\begin{eulercomment}
\eulerheading{2. Tabel}
\begin{eulercomment}
\end{eulercomment}
\eulersubheading{Arti konsep}
\begin{eulercomment}
Data nyata biasanya tersimpan di file berupa tabel (hasil survei, ekspor Excel, dan
sebagainya). EMT dapat membacanya sebagai tabel (frame data), sehingga tidak perlu
diketik ulang. Kolom berisi huruf (misalnya "m" dan "f") dipetakan ke angka, yang disebut
token.

\end{eulercomment}
\eulersubheading{Contoh dunia nyata}
\begin{eulercomment}
Dinas pendidikan menyimpan nilai ribuan siswa di Excel. Dengan membaca file tabel di EMT,
rata-rata dan grafiknya dapat dihitung langsung tanpa mengetik satu per satu.

\end{eulercomment}
\eulersubheading{Praktik di EMT}
\begin{eulercomment}
Pada folder notebook ini ada file table.dat. Isinya hasil survei 25 orang dari buku
berbahasa Jerman "Einführung in die Statistik mit R" karya A. Handl. Empat baris pertama:
\end{eulercomment}
\begin{eulerprompt}
>printfile("table.dat",4);
\end{eulerprompt}
\begin{euleroutput}
  Person Sex Age Titanic Evaluation Tip Problem
  1 m 30 n . 1.80 n
  2 f 23 y g 1.80 n
  3 f 26 y g 1.80 y
\end{euleroutput}
\begin{eulercomment}
Tabel punya 7 kolom: Person, Sex, Age, Titanic, Evaluation, Tip, Problem. Kolom bertanda
huruf perlu diterjemahkan ke angka. Mula-mula saya mendefinisikan himpunan token sendiri.
Fungsi strtokens() mengubah string menjadi vektor token.
\end{eulercomment}
\begin{eulerprompt}
>mf:=["m","f"]; yn:=["y","n"]; ev:=strtokens("g vg m b vb");
\end{eulerprompt}
\begin{eulercomment}
Tabel dibaca dengan terjemahan tersebut. Argumen tok2, tok4, dan seterusnya bukan bagian
dari daftar parameter readtable(), jadi harus diberikan dengan ":=".
\end{eulercomment}
\begin{eulerprompt}
>\{MT,hd\}=readtable("table.dat",tok2:=mf,tok4:=yn,tok5:=ev,tok7:=yn);
>load over statistics;
\end{eulerprompt}
\begin{eulercomment}
Untuk mencetak, himpunan token yang sama harus disebutkan lagi. Di sini hanya 10 baris
pertama yang dicetak.
\end{eulercomment}
\begin{eulerprompt}
>writetable(MT[1:10],labc=hd,wc=5,tok2:=mf,tok4:=yn,tok5:=ev,tok7:=yn);
\end{eulerprompt}
\begin{euleroutput}
   Person  Sex  Age Titanic Evaluation  Tip Problem
        1    m   30       n          .  1.8       n
        2    f   23       y          g  1.8       n
        3    f   26       y          g  1.8       y
        4    m   33       n          .  2.8       n
        5    m   37       n          .  1.8       n
        6    m   28       y          g  2.8       y
        7    f   31       y         vg  2.8       n
        8    m   23       n          .  0.8       n
        9    f   24       y         vg  1.8       y
       10    m   26       n          .  1.8       n
\end{euleroutput}
\begin{eulercomment}
Tanda titik "." menyatakan data yang tidak tersedia (not available).

Jika himpunan token tidak mau ditentukan sendiri, cukup sebutkan kolom mana yang berisi
token, dan readtable() akan mengembalikan himpunan token yang ditemukan.
\end{eulercomment}
\begin{eulerprompt}
>ctok=[2,4,5,7]; \{MT,hd,tok\}=readtable("table.dat",ctok=ctok);
>tok
\end{eulerprompt}
\begin{euleroutput}
  m
  n
  f
  y
  g
  vg
\end{euleroutput}
\begin{eulercomment}
Tabel kini berisi angka hasil terjemahan token. Nilai khusus "." diperlakukan sebagai
NAN (bukan angka). Terjemahan ini dapat diubah dengan parameter NA dan NAval.
\end{eulercomment}
\begin{eulerprompt}
>MT[1]
\end{eulerprompt}
\begin{euleroutput}
  [1,  1,  30,  2,  NAN,  1.8,  2]
\end{euleroutput}
\begin{eulercomment}
Isi tabel dengan angka mentah (belum diterjemahkan kembali):
\end{eulercomment}
\begin{eulerprompt}
>writetable(MT,wc=5)
\end{eulerprompt}
\begin{euleroutput}
      1    1   30    2    .  1.8    2
      2    3   23    4    5  1.8    2
      3    3   26    4    5  1.8    4
      4    1   33    2    .  2.8    2
      5    1   37    2    .  1.8    2
      6    1   28    4    5  2.8    4
      7    3   31    4    6  2.8    2
      8    1   23    2    .  0.8    2
      9    3   24    4    6  1.8    4
     10    1   26    2    .  1.8    2
     11    3   23    4    6  1.8    4
     12    1   32    4    5  1.8    2
     13    1   29    4    6  1.8    4
     14    3   25    4    5  1.8    4
     15    3   31    4    5  0.8    2
     16    1   26    4    5  2.8    2
     17    1   37    2    .  3.8    2
     18    1   38    4    5    .    2
     19    3   29    2    .  3.8    2
     20    3   28    4    6  1.8    2
     21    3   28    4    1  2.8    4
     22    3   28    4    6  1.8    4
     23    3   38    4    5  2.8    2
     24    3   27    4    1  1.8    4
     25    1   27    2    .  2.8    4
\end{euleroutput}
\begin{eulercomment}
Rata-rata satu kolom (misalnya Tip, kolom ke-6) dapat dihitung langsung:
\end{eulercomment}
\begin{eulerprompt}
>mean(tablecol(MT,6))
\end{eulerprompt}
\begin{euleroutput}
  2.175
\end{euleroutput}
\begin{eulercomment}
Rata-rata tip adalah 2,175. Baris yang bertanda "." dilewati oleh tablecol.

Fungsi getstatistics() mengembalikan unsur-unsur berbeda dalam sebuah vektor beserta
banyaknya. Di sini dipakai pada kolom ke-2 (jenis kelamin):
\end{eulercomment}
\begin{eulerprompt}
>\{xu,count\}=getstatistics(tablecol(MT,2)); xu, count,
\end{eulerprompt}
\begin{euleroutput}
  [1,  3]
  [12,  13]
\end{euleroutput}
\begin{eulerprompt}
>writetable(count',labr=tok[xu])
\end{eulerprompt}
\begin{euleroutput}
           m        12
           f        13
\end{euleroutput}
\begin{eulercomment}
Laki-laki (m) 12 orang dan perempuan (f) 13 orang.

Dengan getstatistics() dua kolom, banyaknya setiap kombinasi dapat dihitung:
\end{eulercomment}
\begin{eulerprompt}
>MT24=tablecol(MT,[2,4]); ...
>\{xu1,xu2,count\}=getstatistics(MT24[1],MT24[2]); ...
>writetable(count,labr=tok[xu1],labc=tok[xu2])
\end{eulerprompt}
\begin{euleroutput}
                     n         y
           m         7         5
           f         1        12
\end{euleroutput}
\begin{eulercomment}
Dari 13 perempuan, 12 menjawab y (sudah menonton atau mengenal Titanic). Dari 12
laki-laki, hanya 5 yang menjawab y.

Uji lanjutan (tambahan dari saya): apakah jawaban Titanic bergantung pada jenis kelamin?
H0: kedua hal itu saling bebas. Tabel frekuensi harapan dilihat dulu:
\end{eulercomment}
\begin{eulerprompt}
>writetable(expectedtable(count),wc=8,dc=2,labr=tok[xu1],labc=tok[xu2])
\end{eulerprompt}
\begin{euleroutput}
                 n       y
         m    3.84    8.16
         f    4.16    8.84
\end{euleroutput}
\begin{eulerprompt}
>tabletest(count)
\end{eulerprompt}
\begin{euleroutput}
  0.0066909
\end{euleroutput}
\begin{eulercomment}
Nilai-p sekitar 0,0067 \textless{} 0,05 (dihitung ulang di luar EMT dengan uji chi-kuadrat
Pearson), jadi H0 ditolak. Catatan kritis: dua dari empat sel punya frekuensi harapan
di bawah 5 (3,84 dan 4,16). Pada kondisi itu aproksimasi chi-kuadrat kurang andal, dan
uji eksak Fisher lebih tepat (nilai-p sekitar 0,011, hasilnya tetap menolak H0). Jadi
kesimpulannya sama, tetapi sampel 25 orang membuat dugaan ini perlu dibaca hati-hati.

Tabel dapat ditulis ke berkas, lalu dibaca kembali.
\end{eulercomment}
\begin{eulerprompt}
>filename="test.dat"; ...
>writetable(count,labr=tok[xu1],labc=tok[xu2],file=filename);
>\{MT2,hd2,tk2,hdr\}=readtable(filename,>clabs,>rlabs); ...
>writetable(MT2,labr=hdr,labc=hd2)
\end{eulerprompt}
\begin{euleroutput}
                     n         y
           m         7         5
           f         1        12
\end{euleroutput}
\begin{eulercomment}
Terakhir, berkas dihapus.
\end{eulercomment}
\begin{eulerprompt}
>fileremove(filename);
\end{eulerprompt}
\begin{eulercomment}
\begin{eulercomment}
\eulerheading{3. Distribusi}
\begin{eulercomment}
\end{eulercomment}
\eulersubheading{Arti konsep}
\begin{eulercomment}
Distribusi menggambarkan nilai apa saja yang mungkin muncul dan seberapa sering tiap nilai
muncul. Data hasil simulasi (acak) mengikuti distribusi yang dipilih. Hasilnya berbeda tiap
pengulangan, tetapi polanya tetap. Makin banyak data, makin jelas polanya.

\end{eulercomment}
\eulersubheading{Contoh dunia nyata}
\begin{eulercomment}
Sebelum membuka gerai, pemilik usaha dapat mensimulasikan jumlah pelanggan per jam untuk
menebak kapan antrean paling panjang, tanpa harus mencoba langsung. Simulasi murah dan aman
untuk berlatih sebelum keputusan nyata.

\end{eulercomment}
\eulersubheading{Praktik di EMT}
\begin{eulercomment}
Langkah 1. Dengan plot2d tersedia cara mudah menggambar distribusi data hasil percobaan.
Berikut 1000 data acak berdistribusi normal, histogram dengan 20 kelas, dan kurva normal
baku di atasnya.
\end{eulercomment}
\begin{eulerprompt}
>p=normal(1,1000); //sampel p: 1000 data acak berdistribusi normal
>plot2d(p,distribution=20,style="\(\backslash\)/"); // gambar sampel p
>plot2d("qnormal(x,0,1)",add=1): // tambahkan kurva normal baku
\end{eulerprompt}
\begin{eulercomment}
Perhatikan perbedaan antara batang (sampel) dan kurva (distribusi sebenarnya). Tekan ENTER
pada ketiga perintah itu sekali lagi untuk melihat hasil sampel yang lain.

Langkah 2. Perbandingan 10 simulasi masing-masing 1000 data normal dengan diagram kotak.
Diagram ini menunjukkan median, kuartil 25 \% dan 75 \%, nilai minimum dan maksimum, serta
pencilan.
\end{eulercomment}
\begin{eulerprompt}
>p=normal(10,1000); boxplot(p):
\end{eulerprompt}
\begin{eulercomment}
Langkah 3. Bilangan bulat acak dibangkitkan dengan intrandom. Berikut simulasi pelemparan
dadu 6000 kali. Fungsi getmultiplicities(v,x) menghitung berapa kali unsur v muncul dalam
x, lalu hasilnya digambar dengan columnsplot(). Tiap sisi seharusnya muncul sekitar 1000
kali (garis merah).
\end{eulercomment}
\begin{eulerprompt}
>k=intrandom(1,6000,6);  ...
>columnsplot(getmultiplicities(1:6,k));  ...
>ygrid(1000,color=red):
\end{eulerprompt}
\begin{eulercomment}
Langkah 4. Distribusi eksponensial. Peubah acak kontinu X berdistribusi eksponensial bila
fungsi kerapatannya
\end{eulercomment}
\begin{eulerformula}
\[
f_X(x)=\lambda e^{-\lambda x},\quad x>0,\quad \lambda>0,
\]
\end{eulerformula}
\begin{eulercomment}
dengan parameter
\end{eulercomment}
\begin{eulerformula}
\[
\lambda=\frac{1}{\mu},\quad \mu \text{ adalah rata-rata, ditulis } X \sim \text{Exponential}(\lambda).
\]
\end{eulerformula}
\begin{eulercomment}
Berikut data acak eksponensial dengan rata-rata 2 (cocok untuk waktu tunggu, misalnya
selang waktu antar pelanggan yang datang).
\end{eulercomment}
\begin{eulerprompt}
>plot2d(randexponential(1,1000,2),>distribution):
\end{eulerprompt}
\begin{eulercomment}
Data menumpuk di dekat 0 lalu menurun ke kanan, ciri distribusi eksponensial.

Langkah 5. Untuk banyak distribusi, Euler dapat menghitung fungsi distribusi dan
inversnya. Kurva fungsi distribusi kumulatif normal baku:
\end{eulercomment}
\begin{eulerprompt}
>plot2d("normaldis",-4,4):
\end{eulerprompt}
\begin{eulercomment}
Kurva naik dari 0 ke 1 dan bernilai 0,5 di x = 0.

Langkah 6. Peluang sama dengan luas daerah di bawah kurva kerapatan. Daerah yang diarsir
adalah P(2 \textless{} X \textless{} 5) untuk X normal dengan rata-rata 1 dan simpangan baku 1,5.
\end{eulercomment}
\begin{eulerprompt}
>plot2d("qnormal(x,1,1.5)",-4,6);  ...
>plot2d("qnormal(x,1,1.5)",a=2,b=5,>add,>filled):
>normaldis(5,1,1.5)-normaldis(2,1,1.5)
\end{eulerprompt}
\begin{euleroutput}
  0.24866
\end{euleroutput}
\begin{eulercomment}
Peluangnya 0,2487, yaitu sekitar 24,9 \%. Nilai yang sama diperoleh dari integral
\end{eulercomment}
\begin{eulerformula}
\[
\int_2^5 \frac{1}{1.5\sqrt{2\pi}}e^{-\frac{1}{2}\left(\frac{x-1}{1.5}\right)^2}\,dx
\]
\end{eulerformula}
\begin{eulerprompt}
>gauss("qnormal(x,1,1.5)",2,5)
\end{eulerprompt}
\begin{euleroutput}
  0.24866
\end{euleroutput}
\begin{eulercomment}
\begin{eulercomment}
\eulerheading{4. Distribusi Diskret}
\begin{eulercomment}
\end{eulercomment}
\eulersubheading{Arti konsep}
\begin{eulercomment}
Distribusi diskret memberi peluang pada nilai-nilai yang dapat dihitung (misalnya 1 sampai
6 pada dadu). Euler memungkinkan distribusi diskret buatan sendiri, dan menyediakan
distribusi binomial, hipergeometrik, dan multinomial.

\end{eulercomment}
\eulersubheading{Contoh dunia nyata}
\begin{eulercomment}
Mesin undian yang dicurigai berat sebelah: dengan simulasi dan uji sederhana, kecurangan
sekecil apa pun dapat dideteksi bila percobaannya cukup banyak.

\end{eulercomment}
\eulersubheading{Praktik di EMT}
\begin{eulercomment}
Langkah 1. Distribusi binomial. bindis() memberi peluang paling banyak i keberhasilan dari
n percobaan.
\end{eulercomment}
\begin{eulerprompt}
>bindis(410,1000,0.4)
\end{eulerprompt}
\begin{euleroutput}
  0.7514
\end{euleroutput}
\begin{eulercomment}
Selang kepercayaan Clopper-Pearson untuk parameter p memakai invers fungsi Beta. Artinya,
jika p di luar selang itu, hasil pengamatan 410 dari 1000 termasuk langka.
\end{eulercomment}
\begin{eulerprompt}
>clopperpearson(410,1000)
\end{eulerprompt}
\begin{euleroutput}
  [0.37932,  0.44121]
\end{euleroutput}
\begin{eulercomment}
Cara langsung berikut memberi hasil yang sama, tetapi untuk n besar penjumlahan langsung
tidak teliti dan lambat.
\end{eulercomment}
\begin{eulerprompt}
>p=0.4; i=0:410; n=1000; sum(bin(n,i)*p^i*(1-p)^(n-i))
\end{eulerprompt}
\begin{euleroutput}
  0.7514
\end{euleroutput}
\begin{eulercomment}
invbindis() menghitung invers dari bindis().
\end{eulercomment}
\begin{eulerprompt}
>invbindis(0.75,1000,0.4)
\end{eulerprompt}
\begin{euleroutput}
  409.93
\end{euleroutput}
\begin{eulercomment}
Langkah 2. Pada permainan bridge, ada 5 kartu tersisa (dari 52) pada dua tangan (26
kartu). Peluang pembagian yang lebih buruk daripada 3:2 (misalnya 0:5, 1:4, 4:1, atau 5:0):
\end{eulercomment}
\begin{eulerprompt}
>2*hypergeomsum(1,5,13,26)
\end{eulerprompt}
\begin{euleroutput}
  0.32174
\end{euleroutput}
\begin{eulercomment}
Peluangnya 0,3217, yaitu sekitar 32 \%.

Langkah 3. Simulasi distribusi multinomial: 10 baris, masing-masing 1000 percobaan dengan
peluang 0,4; 0,1; dan 0,5.
\end{eulercomment}
\begin{eulerprompt}
>randmultinomial(10,1000,[0.4,0.1,0.5])
\end{eulerprompt}
\begin{euleroutput}
        390       124       486 
        394        92       514 
        402       135       463 
        399       101       500 
        418        93       489 
        398       102       500 
        393       100       507 
        387        92       521 
        414       102       484 
        413        94       493 
\end{euleroutput}
\begin{eulercomment}
Kolom-kolomnya berkisar di sekitar 400, 100, dan 500.

\begin{eulercomment}
\eulerheading{5. Menggambar Data}
\begin{eulercomment}
\end{eulercomment}
\eulersubheading{Arti konsep}
\begin{eulercomment}
Grafik membuat pola terlihat sekilas. Pilihannya bergantung pada pertanyaan: diagram
batang untuk membandingkan jumlah, diagram lingkaran untuk bagian terhadap keseluruhan,
garis untuk perubahan dari waktu ke waktu, titik (sebar) untuk hubungan dua variabel, dan
histogram untuk bentuk sebaran.

\end{eulercomment}
\eulersubheading{Contoh dunia nyata}
\begin{eulercomment}
Hasil pemilu ditampilkan dengan diagram batang (siapa unggul) dan diagram lingkaran (berapa
persen suara tiap partai). Data yang dipakai adalah hasil pemilu parlemen Jerman sejak 1990
dalam jumlah kursi (dari notebook dosen). Kolom: tahun, jumlah kursi seluruhnya, lalu kursi
CDU/CSU, SPD, FDP, Gr (Hijau), dan Li (Kiri).

\end{eulercomment}
\eulersubheading{Praktik di EMT}
\begin{eulercomment}
\end{eulercomment}
\begin{eulerprompt}
>BW := [ ...
>1990,662,319,239,79,8,17; ...
>1994,672,294,252,47,49,30; ...
>1998,669,245,298,43,47,36; ...
>2002,603,248,251,47,55,2; ...
>2005,614,226,222,61,51,54; ...
>2009,622,239,146,93,68,76; ...
>2013,631,311,193,0,63,64];
\end{eulerprompt}
\begin{eulercomment}
Nama partai disimpan sebagai vektor string:
\end{eulercomment}
\begin{eulerprompt}
>P:=["CDU/CSU","SPD","FDP","Gr","Li"];
\end{eulerprompt}
\begin{eulercomment}
Persentase dicetak rapi. Kolom 3 sampai 7 adalah kursi tiap partai, kolom 2 jumlah seluruh
kursi, dan kolom 1 tahun pemilu.
\end{eulercomment}
\begin{eulerprompt}
>BT:=BW[,3:7]; BT:=BT/sum(BT); YT:=BW[,1]';
\end{eulerprompt}
\begin{eulercomment}
Tabel statistik dicetak dengan nama partai sebagai judul kolom dan tahun sebagai judul
baris. Lebar kolom bawaan wc=10 diperkecil agar lebih padat.
\end{eulercomment}
\begin{eulerprompt}
>writetable(BT*100,wc=6,dc=0,>fixed,labc=P,labr=YT)
\end{eulerprompt}
\begin{euleroutput}
         CDU/CSU   SPD   FDP    Gr    Li
    1990      48    36    12     1     3
    1994      44    38     7     7     4
    1998      37    45     6     7     5
    2002      41    42     8     9     0
    2005      37    36    10     8     9
    2009      38    23    15    11    12
    2013      49    31     0    10    10
\end{euleroutput}
\begin{eulercomment}
Perkalian matriks berikut menjumlahkan persentase dua partai besar. Hasilnya menunjukkan
partai kecil makin berpengaruh sampai 2009.
\end{eulercomment}
\begin{eulerprompt}
>BT1:=(BT.[1;1;0;0;0])'*100
\end{eulerprompt}
\begin{euleroutput}
  [84.29,  81.25,  81.166,  82.753,  72.964,  61.897,  79.873]
\end{euleroutput}
\begin{eulercomment}
Grafik statistik sederhana untuk garis dan titik sekaligus. Alternatifnya, memanggil plot2d
dua kali dengan \textgreater{}add.
\end{eulercomment}
\begin{eulerprompt}
>statplot(YT,BT1,"b"):
\end{eulerprompt}
\begin{eulercomment}
Warna tiap partai ditetapkan lebih dulu:
\end{eulercomment}
\begin{eulerprompt}
>CP:=[rgb(0.5,0.5,0.5),red,yellow,green,rgb(0.8,0,0)];
\end{eulerprompt}
\begin{eulercomment}
Diagram batang perolehan kursi tiap partai pada pemilu 2009:
\end{eulercomment}
\begin{eulerprompt}
>columnsplot(BW[6,3:7],P,color=CP):
\end{eulerprompt}
\begin{eulercomment}
Diagram lingkaran. Karena hitam dan kuning berkoalisi, urutannya disusun ulang.
\end{eulercomment}
\begin{eulerprompt}
>i=[1,3,5,4,2]; piechart(BW[6,3:7][i],color=CP[i],lab=P[i]):
\end{eulerprompt}
\begin{eulercomment}
Diagram bintang menampilkan tiap nilai sebagai sinar dari pusat (data acak):
\end{eulercomment}
\begin{eulerprompt}
>starplot(normal(1,10)+4,lab=1:10,>rays):
\end{eulerprompt}
\begin{eulercomment}
Diagram impuls data acak yang seragam pada [0,1]:
\end{eulercomment}
\begin{eulerprompt}
>plot2d(makeimpulse(1:10,random(1,10)),>bar):
\end{eulerprompt}
\begin{eulercomment}
Diagram pencar (scatter) mudah dibuat dengan plot titik biasa. Grafik berikut menunjukkan
X dan X+Y berkorelasi positif.
\end{eulercomment}
\begin{eulerprompt}
>x=normal(1,100); plot2d(x,x+rotright(x),>points,style=".."):
\end{eulerprompt}
\begin{eulercomment}
\begin{eulercomment}
\eulerheading{6. Regresi dan Korelasi}
\begin{eulercomment}
\end{eulercomment}
\eulersubheading{Arti konsep}
\begin{eulercomment}
Regresi linear mencari garis y = a + b x yang paling dekat dengan titik-titik data. Nilai b
menunjukkan berapa y berubah ketika x naik 1. Korelasi (antara -1 dan 1) mengukur seberapa
erat dua variabel naik atau turun bersama.

\end{eulercomment}
\eulersubheading{Contoh dunia nyata}
\begin{eulercomment}
Jam belajar (x) dan nilai ujian (y): makin lama belajar, nilai cenderung naik. Garis regresi
memberi perkiraan nilai untuk jam belajar tertentu. Perkiraan itu hanya sebatas pola yang
tampak pada data, bukan bukti sebab-akibat.

\end{eulercomment}
\eulersubheading{Praktik di EMT}
\begin{eulercomment}
Langkah 1. Garis regresi untuk data satu peubah dicari dengan polyfit(x,y,1), yang memberi
[a, b].
\end{eulercomment}
\begin{eulerprompt}
>x=1:10; y=[2,3,1,5,6,3,7,8,9,8]; writetable(x'|y',labc=["x","y"])
\end{eulerprompt}
\begin{euleroutput}
           x         y
           1         2
           2         3
           3         1
           4         5
           5         6
           6         3
           7         7
           8         8
           9         9
          10         8
\end{euleroutput}
\begin{eulerprompt}
>p=polyfit(x,y,1)
\end{eulerprompt}
\begin{euleroutput}
  [0.73333,  0.81212]
\end{euleroutput}
\begin{eulercomment}
Garisnya y = 0,733 + 0,812 x.

Langkah 2. Contoh lain: survei mahasiswa berupa usia, usia orang tua, dan banyak saudara
dari file table1.dat. Kolom 2 berisi "m" dan "f"; terjemahannya ditentukan lewat tok2
(bukan dikumpulkan otomatis oleh readtable).
\end{eulercomment}
\begin{eulerprompt}
>\{MS,hd\}:=readtable("table1.dat",tok2:=["m","f"]);  ...
>writetable(MS,labc=hd,tok2:=["m","f"]);
\end{eulerprompt}
\begin{euleroutput}
      Person       Sex       Age    Mother    Father  Siblings
           1         m        29        58        61         1
           2         f        26        53        54         2
           3         m        24        49        55         1
           4         f        25        56        63         3
           5         f        25        49        53         0
           6         f        23        55        55         2
           7         m        23        48        54         2
           8         m        27        56        58         1
           9         m        25        57        59         1
          10         m        24        50        54         1
          11         f        26        61        65         1
          12         m        24        50        52         1
          13         m        29        54        56         1
          14         m        28        48        51         2
          15         f        23        52        52         1
          16         m        24        45        57         1
          17         f        24        59        63         0
          18         f        23        52        55         1
          19         m        24        54        61         2
          20         f        23        54        55         1
\end{euleroutput}
\begin{eulercomment}
Bagaimana usia-usia itu saling bergantung? Kesan pertama didapat dari diagram pencar
berpasangan.
\end{eulercomment}
\begin{eulerprompt}
>scatterplots(tablecol(MS,3:5),hd[3:5]):
\end{eulerprompt}
\begin{eulercomment}
Usia ayah dan usia ibu jelas saling bergantung. Garis regresi dicari: s usia ayah dan t usia
ibu.
\end{eulercomment}
\begin{eulerprompt}
>cs:=MS[,4:5]'; ps:=polyfit(cs[1],cs[2],1)
\end{eulerprompt}
\begin{euleroutput}
  [17.379,  0.74096]
\end{euleroutput}
\begin{eulercomment}
Hasilnya s = 17,38 + 0,741 t. Ini jelas model yang kurang tepat, karena selisih usia ayah
dan ibu mungkin sedikit bergantung pada usia, tetapi tidak sebesar itu. Model yang lebih
wajar adalah s = a + t, dengan a rata-rata selisih usia ayah dan ibu.
\end{eulercomment}
\begin{eulerprompt}
>da:=mean(cs[2]-cs[1])
\end{eulerprompt}
\begin{euleroutput}
  3.65
\end{euleroutput}
\begin{eulercomment}
Rata-rata ayah 3,65 tahun lebih tua daripada ibu. Semua ditampilkan dalam satu diagram
pencar: titik data, garis regresi (merah), dan garis s = t + 3,65 (biru).
\end{eulercomment}
\begin{eulerprompt}
>plot2d(cs[1],cs[2],>points);  ...
>plot2d("evalpoly(x,ps)",color=red,style=".",>add);  ...
>plot2d("x+da",color=blue,>add):
\end{eulerprompt}
\begin{eulercomment}
Diagram kotak kedua usia hanya menunjukkan bahwa keduanya berbeda.
\end{eulercomment}
\begin{eulerprompt}
>boxplot(cs,["mothers","fathers"]):
\end{eulerprompt}
\begin{eulercomment}
Koefisien korelasi menunjukkan korelasi positif.
\end{eulercomment}
\begin{eulerprompt}
>correl(cs[1],cs[2])
\end{eulerprompt}
\begin{euleroutput}
  0.75883
\end{euleroutput}
\begin{eulercomment}
Korelasi peringkat mengukur keserupaan urutan pada kedua vektor, dan juga cukup positif.
\end{eulercomment}
\begin{eulerprompt}
>rankcorrel(cs[1],cs[2])
\end{eulerprompt}
\begin{euleroutput}
  0.75893
\end{euleroutput}
\begin{eulercomment}
Hasilnya 0,759 dan 0,759: hubungan positif yang cukup kuat. Catatan kritis: korelasi kuat
tidak membuktikan bahwa usia satu orang tua menentukan usia yang lain. Model s = t + a
lebih masuk akal dibanding model regresi, walaupun kecocokan angkanya lebih rendah.

\begin{eulercomment}
\eulerheading{7. Membuat Fungsi Baru}
\begin{eulercomment}
Bahasa EMT dapat dipakai untuk memprogram fungsi baru. Sebagai contoh, fungsi kemiringan
(skewness):
\end{eulercomment}
\begin{eulerformula}
\[
\text{sk}(x) = \dfrac{\sqrt{n} \sum_i (x_i-m)^3}{\left(\sum_i (x_i-m)^2\right)^{3/2}}
\]
\end{eulerformula}
\begin{eulercomment}
dengan m rata-rata x. Kemiringan 0 berarti sebaran simetris, positif berarti ekor kanan
lebih panjang, dan negatif berarti ekor kiri lebih panjang.
\end{eulercomment}
\begin{eulerprompt}
>function skew (x:vector) ...
\end{eulerprompt}
\begin{eulerudf}
  m=mean(x);
  return sqrt(cols(x))*sum((x-m)^3)/(sum((x-m)^2))^(3/2);
  endfunction
\end{eulerudf}
\begin{eulercomment}
Bahasa matriks membuat implementasinya sangat singkat dan efisien. Dicoba pada data acak
normal (hasilnya berubah tiap ENTER, sekitar 0 untuk data banyak):
\end{eulercomment}
\begin{eulerprompt}
>data=normal(20); skew(normal(10))
\end{eulerprompt}
\begin{euleroutput}
  -0.13409
\end{euleroutput}
\begin{eulercomment}
Fungsi lain adalah koefisien kemiringan Pearson:
\end{eulercomment}
\begin{eulerprompt}
>function skew1 (x) := 3*(mean(x)-median(x))/dev(x)
>skew1(data)
\end{eulerprompt}
\begin{euleroutput}
  1.423
\end{euleroutput}
\begin{eulercomment}
\begin{eulercomment}
\eulerheading{8. Simulasi Monte Carlo}
\begin{eulercomment}
\end{eulercomment}
\eulersubheading{Arti konsep}
\begin{eulercomment}
Monte Carlo berarti mengulang percobaan acak berkali-kali untuk melihat polanya. Cara ini
dipakai ketika menghitung peluang secara langsung terlalu rumit.

\end{eulercomment}
\eulersubheading{Contoh dunia nyata}
\begin{eulercomment}
Perusahaan asuransi mensimulasikan ribuan skenario klaim untuk memperkirakan dana cadangan
yang diperlukan.

\end{eulercomment}
\eulersubheading{Praktik di EMT}
\begin{eulercomment}
Langkah 1. Simulasi 1000 kali pelemparan 3 dadu, lalu frekuensi tiap jumlah mata (3 sampai
18) dihitung.
\end{eulercomment}
\begin{eulerprompt}
>ds:=sum(intrandom(1000,3,6))';  fs=getmultiplicities(3:18,ds)
\end{eulerprompt}
\begin{euleroutput}
  [4,  19,  27,  41,  72,  100,  118,  118,  138,  108,  77,  75,  57,
  32,  12,  2]
\end{euleroutput}
\begin{eulerprompt}
>columnsplot(fs,lab=3:18):
\end{eulerprompt}
\begin{eulercomment}
Langkah 2. Distribusi yang diharapkan tidak mudah dihitung, sehingga dipakai rekursi.
Fungsi berikut menghitung banyak cara bilangan k dapat dinyatakan sebagai jumlah n bilangan
pada rentang 1 sampai m.
\end{eulercomment}
\begin{eulerprompt}
>function map countways (k; n, m) ...
\end{eulerprompt}
\begin{eulerudf}
    if n==1 then return k>=1 && k<=m
    else
      sum=0;
      loop 1 to m; sum=sum+countways(k-#,n-1,m); end;
      return sum;
    end;
  endfunction
\end{eulerudf}
\begin{eulercomment}
Untuk tiga dadu:
\end{eulercomment}
\begin{eulerprompt}
>cw=countways(3:18,3,6)
\end{eulerprompt}
\begin{euleroutput}
  [1,  3,  6,  10,  15,  21,  25,  27,  27,  25,  21,  15,  10,  6,  3,
  1]
\end{euleroutput}
\begin{eulercomment}
Hasilnya [1, 3, 6, 10, 15, 21, 25, 27, 27, 25, 21, 15, 10, 6, 3, 1], yang berjumlah 216 =
6\textasciicircum{}3. Nilai harapan ditambahkan pada gambar (peluang cw/6\textasciicircum{}3, dikali 1000):
\end{eulercomment}
\begin{eulerprompt}
>plot2d(cw/6^3*1000,>add); plot2d(cw/6^3*1000,>points,>add):
\end{eulerprompt}
\begin{eulercomment}
Batang simulasi mengikuti garis teori. Jumlah beberapa dadu berbentuk lonceng dan mendekati
distribusi normal, yaitu Teorema Limit Pusat: jumlah banyak peubah acak bebas yang identik
mendekati distribusi normal. Contoh nyata: tinggi badan dipengaruhi banyak faktor kecil (gen,
gizi, dan lainnya) yang dijumlahkan, sehingga hasil akhirnya berbentuk lonceng.

Langkah 3. Simpangan baku rata-rata n data normal baku adalah 1/akar(n). Untuk n = 10:
\end{eulercomment}
\begin{eulerprompt}
>longformat; 1/sqrt(10)
\end{eulerprompt}
\begin{euleroutput}
  0.316227766017
\end{euleroutput}
\begin{eulerprompt}
>shortformat;
\end{eulerprompt}
\begin{eulercomment}
Dicek dengan simulasi: 10000 kali vektor acak berisi 10 data.
\end{eulercomment}
\begin{eulerprompt}
>M=normal(10000,10); dev(mean(M)')
\end{eulerprompt}
\begin{euleroutput}
  0.31914
\end{euleroutput}
\begin{eulerprompt}
>plot2d(mean(M)',>distribution):
\end{eulerprompt}
\begin{eulercomment}
Hasil simulasi (sekitar 0,31) mendekati nilai teori 0,316. Artinya, rata-rata banyak data
lebih stabil daripada satu data tunggal. Median dari 10 data normal baku punya simpangan
yang lebih besar:
\end{eulercomment}
\begin{eulerprompt}
>dev(median(M)')
\end{eulerprompt}
\begin{euleroutput}
  0.37273
\end{euleroutput}
\begin{eulercomment}
Hasilnya sekitar 0,365, lebih besar daripada 0,316.

\begin{eulercomment}
\eulerheading{9. Uji Statistika}
\begin{eulercomment}
\end{eulercomment}
\eulersubheading{Arti konsep}
\begin{eulercomment}
Uji statistika menjawab pertanyaan "apakah perbedaan yang terlihat nyata atau hanya
kebetulan". Prosesnya: tetapkan hipotesis nol (H0, biasanya "tidak ada perbedaan"), hitung
nilai-p, lalu putuskan. Semua uji di EMT mengembalikan galat yang diterima jika H0 ditolak,
yaitu nilai-p. Jika nilai-p \textless{} 0,05, H0 ditolak. Jika tidak, H0 tidak dapat ditolak.
"Tidak dapat ditolak" tidak sama dengan "terbukti benar", hanya berarti buktinya belum
cukup.

\end{eulercomment}
\eulersubheading{Contoh dunia nyata}
\begin{eulercomment}
Sebuah sekolah mencoba metode belajar baru dan nilai rata-rata kelas naik 3 poin. Apakah
kenaikan itu karena metodenya atau hanya kebetulan? Uji statistika membantu menjawabnya
dengan angka, bukan dengan perasaan.

\end{eulercomment}
\eulersubheading{Praktik di EMT}
\begin{eulercomment}
a. Uji chi-kuadrat untuk dadu. Pada 600 lemparan didapat frekuensi berikut, yang
dimasukkan ke uji chi-kuadrat. H0: dadu seimbang (seragam).
\end{eulercomment}
\begin{eulerprompt}
>chitest([90,103,114,101,103,89],dup(100,6)')
\end{eulerprompt}
\begin{euleroutput}
  0.49883
\end{euleroutput}
\begin{eulercomment}
Nilai-p 0,4988 \textgreater{} 0,05. Uji chi-kuadrat juga punya mode yang memakai simulasi Monte Carlo.
Parameter \textgreater{}p menganggap vektor y sebagai vektor peluang.
\end{eulercomment}
\begin{eulerprompt}
>chitest([90,103,114,101,103,89],dup(1/6,6)',>p,>montecarlo)
\end{eulerprompt}
\begin{euleroutput}
  0.496
\end{euleroutput}
\begin{eulercomment}
Hasilnya mendekati 0,5 (dosen mendapat 0,529, dan nilainya berubah tiap ENTER karena
simulasi). Galat sebesar itu terlalu besar untuk menolak H0, jadi keseragaman tidak dapat
ditolak. Ini tidak membuktikan dadunya seimbang, hanya bahwa hipotesisnya tidak dapat
ditolak.
b. Uji-t untuk rata-rata. Berikut data dengan rata-rata sebenarnya 200, lalu diuji apakah
rata-ratanya 200. ttest() memerlukan rata-rata, simpangan baku, banyak data, dan nilai
rata-rata yang diuji.
\end{eulercomment}
\begin{eulerprompt}
>s=200+normal([1,100])*10; ...
>ttest(mean(s),dev(s),100,200)
\end{eulerprompt}
\begin{euleroutput}
  0.38342
\end{euleroutput}
\begin{eulercomment}
Karena H0 benar, nilai-p biasanya di atas 0,05 (dosen mendapat 0,2266). Ulangi ENTER dan
nilainya berubah-ubah.

c. Dua pengukuran diuji kesamaan rata-ratanya. H0 ditolak jika hasilnya \textless{} 0,05.
\end{eulercomment}
\begin{eulerprompt}
>tcomparedata(normal(1,10),normal(1,10))
\end{eulerprompt}
\begin{euleroutput}
  0.43805
\end{euleroutput}
\begin{eulercomment}
Jika salah satu distribusi diberi bias (+2), penolakan makin sering. Ulangi beberapa kali
untuk melihat efeknya.
\end{eulercomment}
\begin{eulerprompt}
>tcomparedata(normal(1,10),normal(1,10)+2)
\end{eulerprompt}
\begin{euleroutput}
  0.005791
\end{euleroutput}
\begin{eulercomment}
d. Uji independensi. Dua kelompok orang (laki-laki dan perempuan) memilih satu dari enam
partai.
\end{eulercomment}
\begin{eulerprompt}
>A=[23,37,43,52,64,74;27,39,41,49,63,76];  ...
>  writetable(A,wc=6,labr=["m","f"],labc=1:6)
\end{eulerprompt}
\begin{euleroutput}
             1     2     3     4     5     6
       m    23    37    43    52    64    74
       f    27    39    41    49    63    76
\end{euleroutput}
\begin{eulercomment}
Apakah pilihan bebas dari jenis kelamin? Uji tabel chi-kuadrat menjawabnya.
\end{eulercomment}
\begin{eulerprompt}
>tabletest(A)
\end{eulerprompt}
\begin{euleroutput}
  0.9907
\end{euleroutput}
\begin{eulercomment}
Nilai-p 0,9907 jauh di atas 0,05, jadi kebebasan tidak dapat ditolak. Data ini tidak
menunjukkan bahwa pilihan bergantung pada jenis kelamin. Tabel frekuensi harapan:
\end{eulercomment}
\begin{eulerprompt}
>writetable(expectedtable(A),wc=6,dc=1,labr=["m","f"],labc=1:6)
\end{eulerprompt}
\begin{euleroutput}
             1     2     3     4     5     6
       m  24.9  37.9  41.9  50.3  63.3  74.7
       f  25.1  38.1  42.1  50.7  63.7  75.3
\end{euleroutput}
\begin{eulercomment}
Koefisien kontingensi terkoreksi sangat dekat dengan 0, sehingga pilihan tidak bergantung
pada jenis kelamin.
\end{eulercomment}
\begin{eulerprompt}
>contingency(A)
\end{eulerprompt}
\begin{euleroutput}
  0.042723
\end{euleroutput}
\begin{eulercomment}
e. Analisis ragam (uji-F, ANOVA). Tiga sampel data normal diuji kesamaan rata-ratanya.
Fungsinya varanalysis().
\end{eulercomment}
\begin{eulerprompt}
>x1=[109,111,98,119,91,118,109,99,115,109,94]; mean(x1),
\end{eulerprompt}
\begin{euleroutput}
  106.55
\end{euleroutput}
\begin{eulerprompt}
>x2=[120,124,115,139,114,110,113,120,117]; mean(x2),
\end{eulerprompt}
\begin{euleroutput}
  119.11
\end{euleroutput}
\begin{eulerprompt}
>x3=[120,112,115,110,105,134,105,130,121,111]; mean(x3)
\end{eulerprompt}
\begin{euleroutput}
  116.3
\end{euleroutput}
\begin{eulerprompt}
>varanalysis(x1,x2,x3)
\end{eulerprompt}
\begin{euleroutput}
  0.013805
\end{euleroutput}
\begin{eulercomment}
Rata-ratanya 106,5; 119,1; dan 116,3. Nilai-p 0,0138 \textless{} 0,05, jadi hipotesis kesamaan rata-rata
ditolak, dengan peluang salah 1,4 \%.

f. Dua perlakuan a dan b diterapkan pada orang yang sama. Uji tanda memutuskan apakah a
lebih baik daripada b.
\end{eulercomment}
\begin{eulerprompt}
>a=[8.0,7.4,5.9,9.4,8.6,8.2,7.6,8.1,6.2,8.9];
>b=[6.8,7.1,6.8,8.3,7.9,7.2,7.4,6.8,6.8,8.1];
>signtest(a,b)
\end{eulerprompt}
\begin{euleroutput}
  0.054687
\end{euleroutput}
\begin{eulercomment}
Nilai-p 0,0547 sedikit di atas 0,05, jadi belum dapat ditolak bahwa a sama baiknya dengan b.
Uji Wilcoxon lebih tajam, tetapi memakai besar selisihnya:
\end{eulercomment}
\begin{eulerprompt}
>wilcoxon(a,b)
\end{eulerprompt}
\begin{euleroutput}
  0.029668
\end{euleroutput}
\begin{eulercomment}
Nilai-p 0,0297 \textless{} 0,05, jadi H0 ditolak. Catatan kritis: dua uji yang sama-sama sahih
memberi kesimpulan berbeda pada batas 0,05. Ini tanda bahwa bukti pada data ini hanya
tipis, dan 0,05 bukan "garis ajaib". Dua contoh deret acak berikut menunjukkan perilaku
Wilcoxon: dengan selisih -1 ditolak jelas, tanpa selisih tidak ditolak.
\end{eulercomment}
\begin{eulerprompt}
>wilcoxon(normal(1,20),normal(1,20)-1)
\end{eulerprompt}
\begin{euleroutput}
  0.054213
\end{euleroutput}
\begin{eulerprompt}
>wilcoxon(normal(1,20),normal(1,20))
\end{eulerprompt}
\begin{euleroutput}
  0.11611
\end{euleroutput}
\begin{eulercomment}
\begin{eulercomment}
\eulerheading{10. Catatan untuk Pengguna R}
\begin{eulercomment}
EMT bukan pengganti R sebagai paket statistika, tetapi banyak prosedur statistika tersedia di
EMT. Catatan ini merangkum perbedaan sintaks yang paling sering menjebak bagi yang sudah
mengenal R.\\
- Vektor: operator ":" di EMT dapat memakai ukuran langkah, misalnya 0:0.5:5. Fungsi c() tidak ada, penggabungan memakai kurung siku: [x,0,x].\\
- Indeks negatif: di EMT, x[-2] berarti unsur kedua dari belakang, sedangkan di R berarti "semua kecuali unsur ke-2". Perilaku R ditiru dengan drop(x,2).\\
- Vektor logika: di EMT, hasil perbandingan (0 atau 1) tidak langsung menjadi indeks. Unsur yang bernilai benar diambil dengan nonzeros().
\end{eulercomment}
\begin{eulerprompt}
>x=[10.4, 5.6, 3.1, 6.4, 21.7]; [x,0,x]
\end{eulerprompt}
\begin{euleroutput}
  [10.4,  5.6,  3.1,  6.4,  21.7,  0,  10.4,  5.6,  3.1,  6.4,  21.7]
\end{euleroutput}
\begin{eulerprompt}
>x[-2], drop(x,2)
\end{eulerprompt}
\begin{euleroutput}
  6.4
  [10.4,  3.1,  6.4,  21.7]
\end{euleroutput}
\begin{eulerprompt}
>x>5, x[nonzeros(x>5)]
\end{eulerprompt}
\begin{euleroutput}
  [1,  1,  0,  1,  1]
  [10.4,  5.6,  6.4,  21.7]
\end{euleroutput}
\begin{eulercomment}
\begin{eulercomment}
\eulerheading{11. Membaca dan Menulis Data}
\begin{eulercomment}
\end{eulercomment}
\eulersubheading{Arti konsep}
\begin{eulercomment}
Data biasanya datang dari luar EMT (Excel, survei, web) dan hasil analisis perlu dibawa
keluar. Fungsi writematrix() dan readmatrix() menulis dan membaca matriks. Fungsi
writetable() dan readtable() menulis dan membaca tabel berjudul.

\end{eulercomment}
\eulersubheading{Contoh dunia nyata}
\begin{eulercomment}
Seorang guru mengekspor nilai dari Excel sebagai CSV, menganalisisnya di EMT, lalu menyimpan
ringkasannya kembali sebagai file yang dapat dibuka di Excel.

\end{eulercomment}
\eulersubheading{11.1 Membaca dan menulis vektor}
\begin{eulercomment}
Berikut cara menulis dan membaca vektor bilangan real ke file.
\end{eulercomment}
\begin{eulerprompt}
>a=random(1,100); mean(a), dev(a),
\end{eulerprompt}
\begin{euleroutput}
  0.44364
  0.3037
\end{euleroutput}
\begin{eulercomment}
Data ditulis ke file dengan writematrix(). Untuk catatan sendiri, file ditulis ke folder
yang sama dengan notebook. Vektor kolom a' menghasilkan satu bilangan tiap baris.
\end{eulercomment}
\begin{eulerprompt}
>filename="test.dat";
>writematrix(a',filename);
\end{eulerprompt}
\begin{eulercomment}
Data dibaca dengan readmatrix(), lalu file dihapus.
\end{eulercomment}
\begin{eulerprompt}
>a=readmatrix(filename)';
>fileremove(filename);
>mean(a), dev(a),
\end{eulerprompt}
\begin{euleroutput}
  0.44364
  0.3037
\end{euleroutput}
\begin{eulercomment}
Rata-rata dan simpangan baku sebelum dan sesudah penyimpanan sama, jadi data tersimpan dan
terbaca dengan benar.

writematrix() dan writetable() dapat diatur untuk bahasa lain. Pada sistem Indonesia (koma
sebagai desimal), Excel memerlukan nilai berkoma desimal dengan titik koma sebagai pemisah.
Berkas "test.csv" berikut akan muncul di folder kerja saat ini.
\end{eulercomment}
\begin{eulerprompt}
>filename="test.csv"; ...
>writematrix(random(5,3),file=filename,separator=",");
\end{eulerprompt}
\begin{eulercomment}
Berkas itu dapat langsung dibuka di Excel versi Indonesia, lalu dihapus:
\end{eulercomment}
\begin{eulerprompt}
>fileremove(filename);
\end{eulerprompt}
\begin{eulercomment}
\end{eulercomment}
\eulersubheading{11.2 Berkas CSV}
\begin{eulercomment}
Pertama, sebuah matriks ditulis ke berkas pada folder kerja.
\end{eulercomment}
\begin{eulerprompt}
>file="test.csv";  ...
>M=random(3,3); writematrix(M,file);
>printfile(file)
\end{eulerprompt}
\begin{euleroutput}
  0.1117588409686002,0.05525081823408583,0.2931835879202024
  0.7714979381194919,0.2541355735629442,0.8463894780460306
  0.2938641984846708,0.4657032427307105,0.5283449738795509
  
\end{euleroutput}
\begin{eulercomment}
CSV ini dapat dibuka di Excel versi Inggris dengan klik dua kali. Pada sistem Jerman (atau
Indonesia), data perlu diimpor dengan memperhatikan tanda desimal. Tanda titik desimal juga
bawaan EMT, sehingga matriks dapat dibaca dengan readmatrix().
\end{eulercomment}
\begin{eulerprompt}
>readmatrix(file)
\end{eulerprompt}
\begin{euleroutput}
    0.11176  0.055251   0.29318 
     0.7715   0.25414   0.84639 
    0.29386    0.4657   0.52834 
\end{euleroutput}
\begin{eulercomment}
Di Excel atau program serupa, matriks dapat diekspor sebagai CSV. Contoh: excel-data.csv.
\end{eulercomment}
\begin{eulerprompt}
>printfile("excel-data.csv")
\end{eulerprompt}
\begin{euleroutput}
  0;1000;1000
  1;1051,271096;1072,508181
  2;1105,170918;1150,273799
  3;1161,834243;1233,67806
  4;1221,402758;1323,129812
  5;1284,025417;1419,067549
  6;1349,858808;1521,961556
  7;1419,067549;1632,31622
  8;1491,824698;1750,6725
  9;1568,312185;1877,610579
  10;1648,721271;2013,752707
\end{euleroutput}
\begin{eulercomment}
Sistem Jerman memakai titik koma sebagai pemisah dan koma sebagai desimal. Cara termudah
membacanya ke Euler adalah readmatrix() dengan parameter \textgreater{}comma, yang mengganti koma dengan
titik. Untuk CSV berbahasa Inggris, parameter itu dihilangkan.
\end{eulercomment}
\begin{eulerprompt}
>M=readmatrix("excel-data.csv",>comma)
\end{eulerprompt}
\begin{euleroutput}
          0      1000      1000 
          1    1051.3    1072.5 
          2    1105.2    1150.3 
          3    1161.8    1233.7 
          4    1221.4    1323.1 
          5      1284    1419.1 
          6    1349.9      1522 
          7    1419.1    1632.3 
          8    1491.8    1750.7 
          9    1568.3    1877.6 
         10    1648.7    2013.8 
\end{euleroutput}
\begin{eulercomment}
Data digambar:
\end{eulercomment}
\begin{eulerprompt}
>plot2d(M'[1],M'[2:3],>points,color=[red,green]'):
\end{eulerprompt}
\begin{eulercomment}
Kolom ke-3 (hijau) naik lebih cepat daripada kolom ke-2 (merah). Pola seperti ini sama dengan
dua tabungan yang bunganya berbeda: yang bunganya lebih besar tumbuh lebih cepat. Data ini
terlihat mengikuti pertumbuhan eksponensial: nilai kolom ke-2 pada t=10 adalah 1648,7 =
1000 e\textasciicircum{}0,5 dan kolom ke-3 adalah 2013,8 = 1000 e\textasciicircum{}0,7.

\end{eulercomment}
\eulersubheading{11.3 Membaca dari Web}
\begin{eulercomment}
Situs web atau berkas dengan URL dapat dibuka di EMT dan dibaca baris demi baris. Perlu
koneksi internet. Pada contoh ini, versi terbaru dibaca dari situs EMT dengan mencari
tulisan "Version ..." pada sebuah judul.
\end{eulercomment}
\begin{eulerprompt}
>function readversion () ...
\end{eulerprompt}
\begin{eulerudf}
  urlopen("http://www.euler-math-toolbox.de/Programs/Changes.html");
  repeat
    until urleof();
    s=urlgetline();
    k=strfind(s,"Version ",1);
    if k>0 then substring(s,k,strfind(s,"<",k)-1), break; endif;
  end;
  urlclose();
  endfunction
\end{eulerudf}
\begin{eulerprompt}
>readversion
\end{eulerprompt}
\begin{euleroutput}
  Version 2024-01-12
\end{euleroutput}
\begin{eulercomment}
Hasilnya versi EMT pada halaman tersebut (catatan dosen: Version 2021-04-30). Dapat berbeda
karena situsnya diperbarui. Data publik lain, termasuk data saham, dapat dibaca dengan cara
yang sama, selama dapat diakses.

\begin{eulercomment}
\eulerheading{12. Analisis Data Tambahan dari File yang Dibagikan}
\begin{eulercomment}
Paket dosen menyertakan dua file yang tidak dipakai di notebook dosen: sample.csv dan
GOOG.csv. Keduanya dipakai di sini sebagai latihan analisis data nyata, dengan alur yang
sama: baca, ringkas, gambar, uji.

\end{eulercomment}
\eulersubheading{12.1 Nilai Siswa (sample.csv)}
\begin{eulercomment}
*** Arti konsep dan contoh dunia nyata

File ini berisi 200 baris dengan kolom female (0 = laki-laki, 1 = perempuan), read, write,
math (nilai tes), hon (1 jika masuk kelompok berprestasi), dan femalexmath (hasil kali
female dan math). Contoh nyata: sekolah ingin tahu apakah ada perbedaan nilai antara siswa
laki-laki dan perempuan, dan apakah kemampuan membaca berkaitan dengan kemampuan menulis.

*** Praktik di EMT

Langkah 1. Baca file dengan baris pertama sebagai judul kolom, lalu lihat lima baris pertama.
\end{eulercomment}
\begin{eulerprompt}
>\{SM,hs\}=readtable("sample.csv",>clabs);
>writetable(SM[1:5],labc=hs,wc=12)
\end{eulerprompt}
\begin{euleroutput}
        female        read       write        math         hon femalexmath
             0          57          52          41           0           0
             1          68          59          53           0          53
             0          44          33          54           0           0
             0          63          44          47           0           0
             0          47          52          57           0           0
\end{euleroutput}
\begin{eulercomment}
Langkah 2. Ambil tiap kolom sebagai vektor baris. Nilai asli di file: 109 perempuan dan 91
laki-laki.
\end{eulercomment}
\begin{eulerprompt}
>fm=SM[,1]'; rd=SM[,2]'; wr=SM[,3]'; ma=SM[,4]'; hn=SM[,5]';
>sum(fm), cols(fm)
\end{eulerprompt}
\begin{euleroutput}
  109
  200
\end{euleroutput}
\begin{eulercomment}
Langkah 3. Ringkasan nilai. Hasil hitung ulang di luar EMT: read rata-rata 52,23 (median 50,
simpangan baku 10,25); write 52,78 (54; 9,48); math 52,65 (52; 9,37).
\end{eulercomment}
\begin{eulerprompt}
>mean(rd), median(rd), dev(rd)
\end{eulerprompt}
\begin{euleroutput}
  52.23
  50
  10.253
\end{euleroutput}
\begin{eulerprompt}
>mean(wr), median(wr), dev(wr)
\end{eulerprompt}
\begin{euleroutput}
  52.775
  54
  9.4786
\end{euleroutput}
\begin{eulerprompt}
>mean(ma), median(ma), dev(ma)
\end{eulerprompt}
\begin{euleroutput}
  52.645
  52
  9.3684
\end{euleroutput}
\begin{eulercomment}
Langkah 4. Bandingkan nilai math menurut jenis kelamin. Indeks siswa perempuan dan laki-laki
diambil dengan nonzeros().
\end{eulercomment}
\begin{eulerprompt}
>i=nonzeros(fm==1); j=nonzeros(fm==0);
>mean(ma[i]), mean(ma[j]), dev(ma[i]), dev(ma[j])
\end{eulerprompt}
\begin{euleroutput}
  52.394
  52.945
  9.151
  9.6648
\end{euleroutput}
\begin{eulercomment}
Perempuan 52,39 (simpangan baku 9,15) dan laki-laki 52,95 (9,66). Selisih rata-rata hanya
sekitar 0,55 poin. Uji kesamaan dua rata-rata (H0: sama):
\end{eulercomment}
\begin{eulerprompt}
>tcomparedata(ma[i],ma[j])
\end{eulerprompt}
\begin{euleroutput}
  0.34078
\end{euleroutput}
\begin{eulercomment}
Nilai-p sekitar 0,68 (hasil hitung ulang di luar EMT, 0,68 baik dengan uji Welch maupun
gabungan), jauh di atas 0,05, jadi tidak ada bukti perbedaan nilai math. Untuk nilai write,
perempuan rata-rata 54,99 dan laki-laki 50,12:
\end{eulercomment}
\begin{eulerprompt}
>mean(wr[i]), mean(wr[j])
\end{eulerprompt}
\begin{euleroutput}
  54.991
  50.121
\end{euleroutput}
\begin{eulerprompt}
>tcomparedata(wr[i],wr[j])
\end{eulerprompt}
\begin{euleroutput}
  0.00016982
\end{euleroutput}
\begin{eulercomment}
Nilai-p sekitar 0,0003, jadi H0 ditolak: ada perbedaan nilai menulis, dan perempuan lebih
tinggi sekitar 4,9 poin. Catatan kritis: ini hanya perbedaan rata-rata kelompok pada data
ini, bukan alasan untuk menilai individu, dan sebabnya tidak dapat disimpulkan dari data.

Langkah 5. Hubungan antarnilai. Koefisien korelasi (hitung ulang di luar EMT): read-write
0,597; read-math 0,662; write-math 0,617. Semuanya positif dan cukup kuat.
\end{eulercomment}
\begin{eulerprompt}
>correl(rd,wr), correl(rd,ma), correl(wr,ma)
\end{eulerprompt}
\begin{euleroutput}
  0.59678
  0.66228
  0.61745
\end{euleroutput}
\begin{eulercomment}
Garis regresi write terhadap read:
\end{eulercomment}
\begin{eulerprompt}
>pf=polyfit(rd,wr,1)
\end{eulerprompt}
\begin{euleroutput}
  [23.959,  0.55171]
\end{euleroutput}
\begin{eulercomment}
Hasilnya write = 23,96 + 0,552 read. Setiap kenaikan 1 poin read berkaitan dengan kenaikan
sekitar 0,55 poin write.
\end{eulercomment}
\begin{eulerprompt}
>plot2d(rd,wr,>points); plot2d("evalpoly(x,pf)",>add,color=red):
\end{eulerprompt}
\begin{eulercomment}
Langkah 6. Apakah kelompok berprestasi (hon) bergantung pada jenis kelamin? Tabel
kontingensi dibuat lalu diuji.
\end{eulercomment}
\begin{eulerprompt}
>\{x1,x2,ct\}=getstatistics(fm,hn); ...
>writetable(ct,labr=["laki-laki","perempuan"],labc=["tidak","ya"])
\end{eulerprompt}
\begin{euleroutput}
                 tidak        ya
   laki-laki        74        17
   perempuan        77        32
\end{euleroutput}
\begin{eulerprompt}
>tabletest(ct)
\end{eulerprompt}
\begin{euleroutput}
  0.08043
\end{euleroutput}
\begin{eulercomment}
Frekuensi: laki-laki 74 tidak dan 17 ya (18,7 \%); perempuan 77 tidak dan 32 ya (29,4 \%).
Nilai-p sekitar 0,080 (uji chi-kuadrat Pearson tanpa koreksi, hitung ulang di luar EMT).
Karena 0,080 \textgreater{} 0,05, H0 (saling bebas) tidak dapat ditolak. Selisih 29,4 \% dan 18,7 \% tampak
besar, tetapi pada ukuran sampel ini belum cukup kuat untuk disebut nyata. Catatan kritis:
hasil di dekat 0,05 sebaiknya dilaporkan apa adanya, tanpa dibulatkan ke kesimpulan "ada" atau
"tidak ada".

\end{eulercomment}
\eulersubheading{12.2 Harga Saham Mingguan (GOOG.csv)}
\begin{eulercomment}
*** Arti konsep dan contoh dunia nyata

File ini berisi harga mingguan Google dari 4 November 2019 sampai 2 November 2020 (kolom
Date, Open, High, Low, Close, Adj Close, Volume). Satu baris tanggal 2020-11-02 berisi
"null" (data kosong) dan dibuang, sehingga ada 53 data harga penutupan. Contoh nyata: investor
melihat seberapa besar harga berayun, dan seberapa cepat harga turun pada masa krisis.
Karena kolom tanggal berupa teks dan ada baris "null", harga penutupan (Close) disalin
langsung dari file ke dalam vektor berikut, dalam urutan aslinya.

*** Praktik di EMT

\end{eulercomment}
\begin{eulerprompt}
>Close=[1311.37,1334.87,1295.34,1304.96,1340.62,1347.83,1349.59,1351.89,1360.66, ...
>1429.73,1480.39,1466.71,1434.23,1479.23,1520.74,1485.11,1339.33,1298.41,1219.73, ...
>1072.32,1110.71,1097.88,1211.45,1283.25,1279.31,1320.61,1388.37,1373.19,1410.42, ...
>1428.92,1438.39,1413.18,1431.72,1359.90,1464.70,1541.74,1515.55,1511.87,1482.96, ...
>1494.49,1507.73,1580.42,1644.41,1591.04,1520.72,1459.99,1444.96,1458.42,1515.22, ...
>1573.01,1641.00,1621.01,1638.35];
>n=cols(Close); n
\end{eulerprompt}
\begin{euleroutput}
  53
\end{euleroutput}
\begin{eulercomment}
Ringkasan harga: rata-rata 1414,68; simpangan baku 132,74; terendah 1072,32 (minggu ke-20,
16 Maret 2020, saat pasar jatuh); tertinggi 1644,41 (minggu ke-43, 24 Agustus 2020).
\end{eulercomment}
\begin{eulerprompt}
>mean(Close), dev(Close), min(Close), max(Close)
\end{eulerprompt}
\begin{euleroutput}
  1414.7
  132.74
  1072.3
  1644.4
\end{euleroutput}
\begin{eulerprompt}
>statplot(1:n,Close,"b"):
\end{eulerprompt}
\begin{eulercomment}
Grafik memperlihatkan jatuh tajam pada Maret 2020 lalu pulih. Imbal hasil mingguan dihitung
dari rasio harga berurutan.
\end{eulercomment}
\begin{eulerprompt}
>ret=Close[2:n]/Close[1:n-1]-1;
>mean(ret)*100, dev(ret)*100
\end{eulerprompt}
\begin{euleroutput}
  0.50989
  4.0352
\end{euleroutput}
\begin{eulercomment}
Rata-rata imbal hasil sekitar 0,51 \% per minggu dengan simpangan baku sekitar 4,04 \%. Dari
harga awal 1311,37 sampai akhir 1638,35, kenaikannya sekitar 24,9 \%. Imbal hasil terendah
sekitar -12,1 \% dan tertinggi sekitar 10,3 \% dalam satu minggu.
\end{eulercomment}
\begin{eulerprompt}
>min(ret)*100, max(ret)*100
\end{eulerprompt}
\begin{euleroutput}
  -12.085
  10.344
\end{euleroutput}
\begin{eulercomment}
Garis kecenderungan linear harga terhadap nomor minggu:
\end{eulercomment}
\begin{eulerprompt}
>pg=polyfit(1:n,Close,1)
\end{eulerprompt}
\begin{euleroutput}
  [1274.6,  5.1887]
\end{euleroutput}
\begin{eulercomment}
Hasilnya Close = 1274,6 + 5,19 minggu. Harga naik rata-rata sekitar 5,2 poin tiap minggu.
Catatan kritis: garis lurus ini hanya ringkasan kasar. Penurunan Maret 2020 sangat jauh dari
garis, sehingga tidak layak dipakai untuk meramal. Selain itu, simpangan baku imbal hasil 4 \%
per minggu menunjukkan risiko yang besar. Data hanya satu tahun dan satu saham, sehingga
kesimpulannya tidak dapat digeneralisasi.

\begin{eulercomment}
\eulerheading{13. Latihan Soal}
\begin{eulercomment}
Kelima soal berikut memakai perintah yang sudah dipakai di atas. Jawaban saya hitung ulang di
luar EMT, dan hasilnya sama dengan keluaran yang tertulis di pembahasan. Khusus hasil
simulasi pada Soal 5, angkanya berubah setiap kali perintah dijalankan.

\end{eulercomment}
\eulersubheading{Soal 1: statistika deskriptif}
\begin{eulercomment}
Waktu tempuh (menit) seorang karyawan ke kantor selama 12 hari: 22, 25, 30, 28, 35, 27, 24,
40, 26, 29, 31, 23. Tentukan rata-rata, median, dan simpangan baku, lalu jelaskan apakah ada
nilai yang menarik rata-rata.
\end{eulercomment}
\begin{eulerprompt}
>wt=[22,25,30,28,35,27,24,40,26,29,31,23];
>mean(wt), median(wt), dev(wt)
\end{eulerprompt}
\begin{euleroutput}
  28.333
  27.5
  5.1932
\end{euleroutput}
\begin{eulercomment}
Pembahasan: rata-rata 28,33 menit, median 27,5 menit, simpangan baku 5,19 menit. Rata-rata
sedikit di atas median karena dua hari yang lama (35 dan 40 menit) menariknya ke atas. Median
lebih mewakili hari yang biasa. Simpangan baku 5,19 berarti waktu tempuh biasanya berbeda
sekitar 5 menit dari 28,33.

\end{eulercomment}
\eulersubheading{Soal 2: peluang pada distribusi normal}
\begin{eulercomment}
Berat bersih sebungkus gula berdistribusi normal dengan rata-rata 500 g dan simpangan baku
8 g. Tentukan (a) peluang berat kurang dari 490 g dan (b) peluang berat antara 495 g dan
510 g.
\end{eulercomment}
\begin{eulerprompt}
>normaldis(490,500,8)
\end{eulerprompt}
\begin{euleroutput}
  0.10565
\end{euleroutput}
\begin{eulerprompt}
>normaldis(510,500,8)-normaldis(495,500,8)
\end{eulerprompt}
\begin{euleroutput}
  0.62836
\end{euleroutput}
\begin{eulercomment}
Pembahasan: (a) 0,1056, atau sekitar 10,6 \%. 490 g berada 10/8 = 1,25 simpangan baku di bawah
rata-rata, jadi sekitar 1 dari 10 bungkus lebih ringan dari itu. (b) 0,6284, atau sekitar
62,8 \%. Peluang antara dua nilai adalah selisih fungsi distribusi kumulatif di kedua nilai
itu.

\end{eulercomment}
\eulersubheading{Soal 3: uji chi-kuadrat}
\begin{eulercomment}
Persilangan dua sifat pada kacang polong menghasilkan 556 biji dengan empat jenis: 315, 108,
101, dan 32. Menurut teori Mendel, perbandingannya 9:3:3:1. Ujilah pada taraf 5 \% apakah data
sesuai dengan teori itu.
\end{eulercomment}
\begin{eulerprompt}
>e=556*[9,3,3,1]/16
\end{eulerprompt}
\begin{euleroutput}
  [312.75,  104.25,  104.25,  34.75]
\end{euleroutput}
\begin{eulerprompt}
>chitest([315,108,101,32],e')
\end{eulerprompt}
\begin{euleroutput}
  0.92543
\end{euleroutput}
\begin{eulercomment}
Pembahasan: H0 adalah data sesuai perbandingan 9:3:3:1, sehingga frekuensi harapannya 312,75;
104,25; 104,25; dan 34,75. Statistik chi-kuadrat = 0,016 + 0,135 + 0,101 + 0,218 = 0,470
dengan derajat bebas 3, dan nilai-p sekitar 0,925. Karena 0,925 \textgreater{} 0,05, H0 tidak ditolak:
data tidak bertentangan dengan teori. Ini tidak membuktikan teorinya benar, hanya bahwa data
tidak memberi alasan untuk menolaknya.

\end{eulercomment}
\eulersubheading{Soal 4: regresi linear}
\begin{eulercomment}
Suhu harian (derajat Celsius) dan banyak gelas es teh yang terjual: suhu 24, 26, 28, 30, 32,
34 menghasilkan penjualan 31, 38, 41, 47, 55, 58. Tentukan garis regresi, perkiraan
penjualan pada suhu 29 derajat, dan koefisien korelasinya.
\end{eulercomment}
\begin{eulerprompt}
>sh=[24,26,28,30,32,34]; es=[31,38,41,47,55,58];
>ps=polyfit(sh,es,1)
\end{eulerprompt}
\begin{euleroutput}
  [-34.543,  2.7429]
\end{euleroutput}
\begin{eulerprompt}
>evalpoly(29,ps), correl(sh,es)
\end{eulerprompt}
\begin{euleroutput}
  45
  0.99307
\end{euleroutput}
\begin{eulerprompt}
>statplot(sh,es,"b"); plot2d("evalpoly(x,ps)",>add,color=blue,style="--"):
\end{eulerprompt}
\begin{eulercomment}
Pembahasan: garisnya y = -34,54 + 2,743 x. Setiap kenaikan suhu 1 derajat menambah penjualan
sekitar 2,7 gelas. Untuk suhu 29 derajat, perkiraannya 45 gelas. Korelasi 0,993 menunjukkan
hubungan positif yang sangat kuat. Suhu 29 berada di dalam rentang data (24 sampai 34), jadi
perkiraannya lebih dapat dipercaya daripada perkiraan di luar rentang. Konstanta -34,54 tidak
punya arti nyata, karena suhu 0 derajat jauh di luar data.

\end{eulercomment}
\eulersubheading{Soal 5: binomial dan simulasi}
\begin{eulercomment}
Sebuah kuis berisi 10 soal pilihan ganda dengan 4 pilihan, dan seorang siswa menjawab dengan
menebak. Hitung (a) peluang tepat 3 jawaban benar dan (b) peluang paling sedikit 6 jawaban
benar. Cek dengan simulasi 10000 siswa.
\end{eulercomment}
\begin{eulerprompt}
>bindis(3,10,0.25)-bindis(2,10,0.25)
\end{eulerprompt}
\begin{euleroutput}
  0.25028
\end{euleroutput}
\begin{eulerprompt}
>1-bindis(5,10,0.25)
\end{eulerprompt}
\begin{euleroutput}
  0.019728
\end{euleroutput}
\begin{eulerprompt}
>kb=sum(intrandom(10000,10,4)==1)'; mean(kb==3), mean(kb>=6)
\end{eulerprompt}
\begin{euleroutput}
  0.2537
  0.0179
\end{euleroutput}
\begin{eulercomment}
Pembahasan: banyak jawaban benar berdistribusi binomial dengan n = 10 dan p = 0,25.
(a) Peluang tepat 3 benar adalah selisih bindis(3) dan bindis(2), yaitu 0,2503. (b) Peluang
paling sedikit 6 benar adalah 1 - bindis(5) = 0,0197, atau sekitar 2 \%, jadi menebak hampir
tidak pernah menghasilkan 6 benar atau lebih. Hasil simulasi berada di sekitar 0,25 dan 0,02
dan berubah di setiap ENTER. Makin banyak siswa yang disimulasikan, hasilnya makin dekat ke
nilai teori.

\begin{eulercomment}
\eulerheading{Kesimpulan}
\begin{eulercomment}
1. Data statistik disimpan sebagai vektor, matriks, dan tabel, serta dibaca dari file.
2. Statistika deskriptif (rata-rata, median, simpangan baku, kuartil) meringkas data.
3. Data acak membantu melihat pola distribusi. Pola itu makin jelas pada data banyak.
4. Uji statistika memakai nilai-p: H0 ditolak jika nilai-p \textless{} 0,05. Nilai-p dekat 0,05 perlu
\end{eulercomment}
\begin{eulerttcomment}
   dibaca hati-hati, dan "tidak ditolak" bukan berarti "terbukti benar".
\end{eulerttcomment}
\begin{eulercomment}
5. Simulasi jumlah dadu menunjukkan Teorema Limit Pusat: hasilnya mendekati normal.
6. Regresi dan korelasi menunjukkan hubungan antardata, tetapi tidak membuktikan sebab-akibat.
7. Pada data nyata (sample.csv dan GOOG.csv), perbedaan rata-rata, korelasi, dan risiko dapat
\end{eulercomment}
\begin{eulerttcomment}
   dihitung dengan perintah yang sama, tetapi kesimpulannya harus mempertimbangkan ukuran
   sampel dan asumsi uji.
\end{eulerttcomment}
\end{eulernotebook}
\end{document}
